\documentclass[10pt,a4paper]{amsart}
\usepackage[margin=19mm]{geometry}
\usepackage{amsmath,amssymb,mathtools}
\usepackage[scr=rsfs,cal=euler]{mathalfa}
\usepackage{xfrac}
\usepackage[unicode,hidelinks,bookmarks=false]{hyperref}
\newcommand{\ie}{i.e.\ }
\newcommand{\cf}{cf.\ }
\newcommand{\resp}{respectively\ }
\newcommand{\Subsection}{\subsection}
\providecommand{\nobreakdash}{\nobreak\hbox{-}}
\setlength{\emergencystretch}{2em}
\multlinegap0pt
\usepackage{amssymb}
\usepackage[shortlabels]{enumitem}
\newtheorem{lemma}{Lemma}
\DeclareMathOperator{\supp}{supp}
\title{Crossed product $C^*$-algebras associated with non-minimal actions on the circle}
\author{Jamie Bell}
\date{Source: arXiv:2604.18410v2 (CC BY 4.0)}
\begin{document}
\maketitle

\noindent\emph{Source and licence.} Crossed product $C^*$-algebras associated with non-minimal actions on the circle. Version arXiv:2604.18410v2 (CC BY 4.0), distributed by arXiv under the Creative Commons Attribution 4.0 International licence, http://creativecommons.org/licenses/by/4.0/, as recorded in the arXiv licence metadata for this exact version and shown on the abstract page https://arxiv.org/abs/2604.18410v2. Full licence notice: \texttt{licenses/arxiv-2604.18410-CC-BY-4.0.txt}.\par\smallskip

\noindent\textit{Editorial scope.} Counted unit: the article's lemma supp (source0.tex lines 255--257). The article's complete original proof of it is delivered with it (source0.tex lines 259--261, one \texttt{proof} environment of 3 lines). No mathematics is written, repaired or re-derived: the statement and the proof are the article's own lines, in the article's own notation, and no retained line is substituted or re-flowed.  No further source-local context is needed: every cross-reference of every retained line resolves inside the two delivered spans.  Nothing was omitted from any retained span, so the package carries no omission record.  No retained line contains a citation, so the wrapper carries no bibliography and no bibliography entry.  No retained line contains a figure environment, an image-inclusion command or drawing code, so no image is delivered and the package declares no asset dependency.  Editorial formatting only, in addition: an AMS \texttt{amsart} preamble that restates the article's own declarations and macros used by the retained lines, namely the environments \texttt{lemma} on separate counters, together with the standard packages those lines need.  The retained environments print in this document as Lemma 1 in order of appearance, whereas the article prints the same items with the numbering of its own full document.  The complete native source of the article as distributed in the arXiv e-print for this exact version is bundled unchanged as \texttt{sources/198797/source0.tex}.
\begin{lemma}\label{L-minimal supp}
    Let $G\curvearrowright X$ be an action of a discrete amenable group on a compact Hausdorff space $X$ and suppose that $M_G(X) = \{\mu\}$. Then $\supp(\mu)$ is the unique minimal set. In particular, $G\curvearrowright X$ is minimal if and only if $\mu$ has full support.  
\end{lemma}

\begin{proof}
     Since $\supp(\mu)$ is a non-empty closed (hence compact) and $G$-invariant set, it contains a minimal set $Y\subseteq \supp(\mu)$. Since $G$ is amenable, there exists a Borel measure $\nu \in M_G(Y)\subseteq M_G(X)$ such that $\nu(Y) = 1$. By unique ergodicity, we have $\nu = \mu$ and hence $\mu(Y) = 1$. It follows from the definition of $\supp(\mu)$ that $\supp(\mu)\subseteq Y$, hence $\supp(\mu) = Y$ is a minimal set. If $Y_1,Y_2\subseteq X$ are two minimal sets then since $Y_1\cap Y_2$ is also closed and $G$-invariant, either $Y_1 = Y_2$ or $Y_1$ and $Y_2$ are disjoint. If $Y_1$ and $Y_2$ are disjoint, then $\mu_i := \mu|_{Y_i}$ determine two distinct probability measures, contradicting uniqueness of $\mu$. Hence $\supp(\mu)$ is the unique minimal set. 
\end{proof}

\end{document}
